% s-063-p063-m002; original PDF 63, 64
\Aphorism{s-063-p063-m002}{62}
Jede Synthese ist Gewinn, jede Analyse Verlust. Eine Zusammensetzung, die nicht mehr als Summe ihrer Theile erkennbar ist, wiegt einer neuen Realität gleich; ein Auseinandernehmen vermehrt die Anzahl der niederen, eigenschaftsloseren Complexe und verringert die der höheren, qualificirteren. Einem Gehirn, das zwei Begriffe hat, gebe ich einen dritten, indem ich zwischen beiden einen \Foreign{nexus} suche, der so gesucht und willkürlich wie möglich sein darf; der Zwang, zwei Unverträgliche zusammenzudenken, löthet endlich beide wirklich zusammen. Also Ehen schliessen, der Kindererzeugung wegen. Einen Menschen, der fertig ist, bereichere ich ethisch, indem ich ihn in unmögliche Situationen bringe, die er mit seiner bisherigen Ethik nicht verdaut; der Nothbehelf, den er ersinnt und den ich ihm durch Wiederholung einübe, ist dann ein neues Hülfsmittel, ein Zuwachs an Lebenskenntniss. So hat der Mann von Ehre vor dem gewöhnlichen Menschen den Vorzug, in ein paar schwierige und unmögliche Lagen mehr kommen zu können, als Jener, also verschiedene Male mehr zur Aufbietung aller Kraft und Erfindsamkeit genöthigt zu sein. Der Ehrbegriff selbst ist nichts anderes als der Zwang, alle Dinge von einer Seite mehr zu betrachten als alle Welt; er bereichert also das Aussehen der Dinge, er vervielfacht unsere Interessen und Beziehungen, er vertritt als technische Verlängerung die Stelle eines fehlenden oder noch unentwickelten Organs und trägt zugleich zu dessen Züchtung bei. Jede moralische oder ästhetische Idiosynkrasie, die wir uns aneignen, verschönt die Welt (ohne uns freilich immer das Leben bequem zu machen); umgekehrt wirkt jeder Indifferentismus dahin, sie zu verwischen, die Farben auszulöschen und uns den Quell förderlicher und bildender Erlebnisse zuzuschütten. Man soll, bei und in sich, alle Arten von Tyrannei und Vergewaltigung unterstützen, alle Arten von \Foreign{laisser aller} und gutmüthiger Erleichterung in die Acht erklären. Man lasse sich nicht gehen, man exercire sich lieber, um gehen zu lernen. Man glaube dem Augenschein nicht, der diese Disciplin für unnöthig hält; man mache sich das Leben schwer, weil alle Welt es sich leicht macht. Man verspreche, um halten zu müssen; man verlange von sich Alles, um sicher zu sein, nichts vergeblich zu verlangen; man erfinde fixe Ideen, wenn die überlieferten nicht genügen. Man begebe sich in zehn Gefahren mehr als der Nachbar, unter zehn Möglichkeiten des Missverständnisses mehr, in zehn Versuchungen mehr; man beurtheile sich nach einem härteren Gesetz und mit schärferer Betonung der erschwerenden Umstände, man vervielfache seine Casuistik, nicht um öfter davonzukommen, sondern um öfter zu straucheln. Man verschmälere seine Lebensgrundlage, um öfter das Gleichgewicht zu verlieren; man entferne sein Ziel, um leichter den Weg zu verfehlen; man springe hinüber, wo alle Welt überfährt, und verbinde sich die Augen, wo die Anderen sie bewaffnen. Man habe einen Gott mehr, eine Moral mehr, eine Ehre mehr als Jedermann, zu keinem anderen Zwecke, als um mehr zu erleben als Jedermann.
\par

